% ===============================================================================
% LH-06b: Lehrerhinweise - ir + estar (Bewegung und Lage)
% Abgeleitet aus LE-06b (Verlauf, Differenzierung, Fehler)
% ===============================================================================

\documentclass[11pt, a4paper]{article}

\usepackage[utf8]{inputenc}
\usepackage[T1]{fontenc}
\usepackage[ngerman]{babel}
\usepackage{helvet}
\renewcommand{\familydefault}{\sfdefault}

\usepackage[margin=2cm]{geometry}
\usepackage{enumitem}
\usepackage{tabularx}
\usepackage{booktabs}
\usepackage{longtable}

\usepackage{xcolor}
\usepackage{tcolorbox}

\usepackage{fancyhdr}
\usepackage{lastpage}
\usepackage{amssymb}

% Farben
\definecolor{spanisch}{HTML}{c41e3a}
\definecolor{grau}{HTML}{888888}
\definecolor{hellgrau}{HTML}{f5f5f5}
\definecolor{basis}{HTML}{2a7a4b}
\definecolor{standard}{HTML}{2c5aa0}
\definecolor{erweiterung}{HTML}{7b1fa2}
\definecolor{loesung}{HTML}{c62828}

\newcommand{\fachfarbe}{spanisch}

% Variablen
\newcommand{\thema}{ir + estar (Bewegung und Lage)}
\newcommand{\sequenz}{06}
\newcommand{\block}{b}
\newcommand{\fach}{Spanisch}
\newcommand{\klasse}{7}
\newcommand{\dauer}{90 Min. (Doppelstunde)}

% Kopf-/Fußzeile
\pagestyle{fancy}
\fancyhf{}
\renewcommand{\headrulewidth}{0.4pt}
\renewcommand{\footrulewidth}{0.4pt}
\lhead{\fach{} -- Klasse \klasse}
\chead{\textbf{LH-\sequenz\block{} (Lehrerhinweise)}}
\rhead{}
\lfoot{}
\rfoot{Seite \thepage\ von \pageref{LastPage}}

% Differenzierungsmarker
\newcommand{\B}{\textcolor{basis}{\textbf{[B]}}}
\newcommand{\Std}{\textcolor{standard}{\textbf{[S]}}}
\newcommand{\E}{\textcolor{erweiterung}{\textbf{[E]}}}
\newcommand{\loesung}[1]{\textcolor{loesung}{\textbf{#1}}}
\newcommand{\es}[1]{\textcolor{\fachfarbe}{\textbf{#1}}}

% Boxen
\newtcolorbox{kurzplan}{
    colback=hellgrau,
    colframe=\fachfarbe,
    boxrule=1pt,
    arc=0pt,
    fonttitle=\bfseries,
    title=Kurzplan
}

\newtcolorbox{differenzierung}{
    colback=white,
    colframe=grau,
    boxrule=0.5pt,
    arc=0pt,
    fonttitle=\bfseries,
    title=Differenzierung
}

\newtcolorbox{fehlerbox}{
    colback=white,
    colframe=loesung,
    boxrule=0.5pt,
    arc=0pt,
    fonttitle=\bfseries,
    title=Häufige Fehler
}

\newtcolorbox{materialbox}{
    colback=white,
    colframe=\fachfarbe,
    boxrule=0.5pt,
    arc=0pt,
    fonttitle=\bfseries,
    title=Material-Checkliste
}

% ===============================================================================
\begin{document}

\begin{center}
    {\Large\bfseries\textcolor{\fachfarbe}{Lehrerhinweise: \thema}}\\[0.3em]
    {\normalsize LH-\sequenz\block{} | \dauer{} | Std. 3-4 von 8}
\end{center}

\vspace{10pt}

% ===============================================================================
% 1. KURZPLAN
% ===============================================================================

\section*{1. Stundenverlauf (Kurzplan)}

\textbf{1. Stunde (45 Min.) -- Fokus: ir + a}

\begin{tabularx}{\textwidth}{|c|l|X|l|}
\hline
\textbf{Zeit} & \textbf{Phase} & \textbf{Inhalt} & \textbf{Material} \\
\hline
0--5 & Einstieg & Stadtplan: ¿Adónde vas los sábados? & Stadtplan \\
\hline
5--10 & Aktivierung & Orte wiederholen, Beispiel: Voy al cine & Tafel \\
\hline
10--20 & Input ir & Verb \textit{ir} einführen, Kontraktion \textit{al} & Tafel, AB \\
\hline
20--30 & Übung 1 & Chorisches Sprechen, Konjugation & -- \\
\hline
30--40 & Übung 2 & PA: Minidialoge ¿Adónde vas? & Stadtplan \\
\hline
40--45 & Sicherung & Merkkasten ir ausfüllen & AB-06b \\
\hline
\end{tabularx}

\vspace{1em}

\textbf{2. Stunde (45 Min.) -- Fokus: estar}

\begin{tabularx}{\textwidth}{|c|l|X|l|}
\hline
\textbf{Zeit} & \textbf{Phase} & \textbf{Inhalt} & \textbf{Material} \\
\hline
0--5 & Aktivierung & Quiz: ir-Formen abfragen & -- \\
\hline
5--15 & Input estar & Verb \textit{estar} für Ortsangaben & Tafel, AB \\
\hline
15--25 & Übung 3 & Stadtplan: ¿Dónde está el cine? & Stadtplan \\
\hline
25--35 & Anwendung & Kombinationsübung: ir + estar & AB-06b \\
\hline
35--40 & Sicherung & Merkkasten estar, Unterschied ir/estar & AB-06b \\
\hline
40--45 & Abschluss & Zusammenfassung, HA & -- \\
\hline
\end{tabularx}

% ===============================================================================
% 2. DIFFERENZIERUNG
% ===============================================================================

\section*{2. Differenzierung}

\begin{differenzierung}
\textbf{Legende:} \B{} Basis (BR) | \Std{} Standard (MR) | \E{} Erweiterung (GY)

\vspace{0.5em}
\textbf{WICHTIG:} Diese Marker sind NUR für die Lehrkraft!
\end{differenzierung}

\vspace{1em}

\begin{tabularx}{\textwidth}{|l|X|X|X|}
\hline
& \B{} \textbf{Basis} & \Std{} \textbf{Standard} & \E{} \textbf{Erweiterung} \\
\hline
\textbf{Aufg. 1} & Wortbank mit allen Formen & Formen selbst einsetzen & + Sätze verlängern \\
\hline
\textbf{Aufg. 2} & Regelhilfe sichtbar & Regel anwenden & Ohne Hilfe + Begründung \\
\hline
\textbf{Aufg. 3} & Lücken mit Auswahl & Frei einsetzen & + eigene Sätze bilden \\
\hline
\textbf{Aufg. 4} & Dialog mit Vorlage (4 Zeilen) & Vollständig (6 Zeilen) & + Stadtbeschreibung \\
\hline
\end{tabularx}

\vspace{1em}

\textbf{Konkrete Hinweise:}
\begin{itemize}
    \item \B{} LRS-SuS: Konjugationstabelle als Hilfe erlauben, Zeitzugabe (+10 Min.)
    \item \B{} DaZ-SuS: Ortsnamen deutsch-spanisch vorab besprechen
    \item \Std{} Standardgruppe: Wechselnde Partner bei Stadtplan-Dialog
    \item \E{} Leistungsstarke: Stadtplan selbst beschriften, Wegbeschreibung
\end{itemize}

% ===============================================================================
% 3. HÄUFIGE FEHLER
% ===============================================================================

\section*{3. Häufige Fehler}

\begin{fehlerbox}
\begin{tabularx}{\textwidth}{|l|X|X|}
\hline
\textbf{Fehler} & \textbf{Beispiel} & \textbf{Reaktion} \\
\hline
Kontraktion vergessen & *``Voy a el cine'' & Regel visualisieren: a + el = al \\
\hline
ser statt estar & *``El cine es en la plaza'' & Merksatz: estar = wo, ser = wer \\
\hline
Akzent vergessen & *``esta'' statt ``está'' & Betonung üben: estÁ (betont!) \\
\hline
Falsche Person & *``Yo va al cine'' & Person-Endung-Zuordnung üben \\
\hline
\end{tabularx}
\end{fehlerbox}

% ===============================================================================
% 4. MATERIAL-CHECKLISTE
% ===============================================================================

\section*{4. Material-Checkliste}

\begin{materialbox}
\begin{itemize}[label=$\square$]
    \item AB-06b.pdf (24 Kopien)
    \item ML-06b.pdf (1x für Lehrkraft)
    \item Lehrwerk: Qué pasa 1, Unidad 6
    \item Stadtplan (Kopiervorlage, 12 Sets für PA)
    \item Optional: LP-06b.html (Tablets/PCs)
    \item Tafel/Whiteboard: Konjugationstabellen
\end{itemize}
\end{materialbox}

% ===============================================================================
% 5. LÖSUNGEN
% ===============================================================================

\section*{5. Lösungen AB-06b}

\textbf{Merkkasten ir:}
voy, vas, va, vamos, vais, van

\textbf{Merkkasten estar:}
estoy, estás, está, estamos, estáis, están

\vspace{0.5em}

\textbf{Aufgabe 1:} (je 1P)
\begin{enumerate}[label=\alph*)]
    \item \loesung{voy}
    \item \loesung{vas}
    \item \loesung{va}
    \item \loesung{vamos}
\end{enumerate}

\textbf{Aufgabe 2:} (je 1P)
\begin{enumerate}[label=\alph*)]
    \item \loesung{al} (a + el = al)
    \item \loesung{a la} (la $\rightarrow$ keine Kontraktion)
    \item \loesung{al} (a + el = al)
    \item \loesung{a la} (la $\rightarrow$ keine Kontraktion)
\end{enumerate}

\textbf{Aufgabe 3:} (je 1,5P)
\begin{enumerate}[label=\alph*)]
    \item El cine \loesung{está} en la plaza.
    \item La biblioteca \loesung{está} cerca del parque.
    \item ¿Dónde \loesung{estás}? -- \loesung{Estoy} en casa.
\end{enumerate}

\textbf{Aufgabe 4:} (6P)
\begin{itemize}
    \item ir-Form korrekt: 2P
    \item estar-Form korrekt: 2P
    \item Sinnvolle Struktur: 1P
    \item Sprachliche Korrektheit: 1P
\end{itemize}

% ===============================================================================
% 6. HAUSAUFGABE
% ===============================================================================

\section*{6. Hausaufgabe}

\begin{itemize}
    \item Lehrwerk Unidad 6, Übung zu ir/estar
    \item Vokabeln lernen: ir-Formen, estar-Formen, Orte
    \item Optional: Lernpfad LP-06b.html bearbeiten
    \item \E{}: Kurzer Text ``Mi ciudad'' (5 Sätze mit ir/estar)
\end{itemize}

% ===============================================================================
% 7. REFLEXIONSNOTIZEN
% ===============================================================================

\section*{7. Reflexionsnotizen (nach Durchführung)}

\begin{tabularx}{\textwidth}{|l|X|}
\hline
\textbf{Was lief gut?} & \\[2em]
\hline
\textbf{Was lief nicht?} & \\[2em]
\hline
\textbf{Zeitmanagement} & \\[2em]
\hline
\textbf{Für nächstes Mal} & \\[2em]
\hline
\end{tabularx}

\end{document}
